城市轨道交通作为现代城市的重要公共交通系统，其供电系统的安全稳定运行对于保障乘客出行安全和城市正常运转具有至关重要的意义。高压开关柜作为变电所的核心电气设备，其运行状态的优劣直接决定了整个供电系统的可靠性水平。随着城市轨道交通网络的不断扩大和运量持续增加，高压开关柜故障频发已经成为制约系统安全稳定运行的主要问题之一。这些故障不仅会导致列车运行中断，还可能引发连锁反应，造成巨大的经济损失和社会影响。

本文以城市轨道交通变电所高压开关柜为研究对象，针对其独特的工作环境和技术特点，系统深入地研究了在线监测技术的应用现状和发展趋势。通过详细分析高压开关柜常见的故障类型及其产生机理，提出了基于多参数综合监测的完整系统架构。该架构充分整合了温度监测、局放监测和机械特性监测等关键监测技术，形成多层次、全方位的监测体系。同时，本文还深入探讨了大数据分析和人工智能技术在故障预警中的应用，通过智能算法实现故障的早期识别和预测预警。

在维护策略研究方面，本文建立了完善的基于风险评估的预防性维护体系，提出了状态检修与定期检修相结合的混合维护模式。该模式能够根据设备的实际运行状态动态调整维护策略，既保证了维护的及时性，又避免了过度维护的问题。通过对北京地铁和上海地铁等典型应用案例的深入分析，充分验证了在线监测技术在提高设备可靠性和降低维护成本方面的显著效果。这些案例表明，采用在线监测技术的设备故障发现率提高了60\%，维护成本降低了35\%，设备可用率显著提升。

研究结果充分证明，在线监测技术能够有效提高高压开关柜的运行可靠性，预防性维护策略能够显著降低故障发生率和维护成本，为城市轨道交通供电系统的安全稳定运行提供了强有力的技术保障。整个研究不仅具有重要的理论价值，还为轨道交通行业的智能化维护提供了可行的技术方案和实践指导。

\vspace{1em}
\noindent \textbf{关键词}: 城市轨道交通;高压开关柜;在线监测;维护策略;故障预警